\documentclass[12pt]{article}
\usepackage[T1]{fontenc}
\usepackage[utf8]{inputenc}
\usepackage{lmodern}
\usepackage{textcomp}
\usepackage{enumitem}
\usepackage{geometry}
\geometry{margin=1in}
\begin{document}
\begin{center}
Answer Key - Health Sciences - Graduate Level Assessment 5 \
\small (Answers are ordered exactly as in the question paper.)
\end{center}

\section*{Multiple Choice}
\begin{enumerate}[label=\arabic*.]
\item \textbf{Answer:} A
\\ \textit{Options:} A) Gerodiversity; B) Geropsychiatry; C) Geronutrition; D) Gerolinguistics; E) Gerocommunication
\\ \textit{Note:} The term "gerodiversity" accurately reflects the study of the diverse experiences, needs, and characteristics of older adults, emphasizing the complexity of aging within health sciences.
\item \textbf{Answer:} D
\\ \textit{Options:} A) India, Brazil, and South Africa; B) Canada, Mexico, and Japan; C) Australia, New Zealand, and Indonesia; D) China, the U.S. and Russia; E) Afghanistan, Iraq and Myanmar
\\ \textit{Note:} The correct answer is based on research data indicating that China, the U.S., and Russia have the highest reported numbers of individuals who engage in injectable drug use, reflecting significant public health challenges related to substance abuse in these countries.
\item \textbf{Answer:} A
\\ \textit{Options:} A) involved in the synthesis of ATP; B) types of RNA-binding proteins; C) part of the hemoglobin in blood cells; D) involved in the control of translation; E) components of the cell's cytoskeleton
\\ \textit{Note:} Zinc finger proteins and helix-turn-helix proteins play crucial roles in gene regulation, which is essential for the expression of enzymes involved in ATP synthesis.
\item \textbf{Answer:} A
\\ \textit{Options:} A) Uric acid; B) Nitrous Oxide; C) Nitric Acid; D) Nitrites; E) Nitrogen
\\ \textit{Note:} Uric acid is the main nitrogenous compound in urine because it is the primary end product of purine metabolism in humans, reflecting the body's method of excreting nitrogen waste.
\item \textbf{Answer:} C
\\ \textit{Options:} A) protein.; B) cytosine.; C) thymine.; D) guanine.; E) nucleotide.
\\ \textit{Note:} Adenine pairs with thymine in DNA due to the complementary base pairing rules established by hydrogen bonding, where adenine forms two hydrogen bonds with thymine, ensuring accurate DNA replication and stability.
\end{enumerate}

\section*{True / False}
\begin{enumerate}[label=\arabic*.]
\setcounter{enumi}{5}
\item \textbf{Answer:} True
\\ \textit{Options:} A) True; B) False
\\ \textit{Note:} Public health ethics prioritizes the well-being of entire populations by addressing social determinants of health, health disparities, and collective interventions, distinguishing it from individual-centered medical ethics.
\item \textbf{Answer:} False
\\ \textit{Options:} A) True; B) False
\\ \textit{Note:} Angelman syndrome is primarily caused by a deletion or mutation in the UBE 3 A gene on chromosome 15, rather than a mutation in the SNRPN gene, which is involved in genomic imprinting but not the direct cause of the syndrome.
\item \textbf{Answer:} False
\\ \textit{Options:} A) True; B) False
\\ \textit{Note:} An ileostomy opening is typically larger than 1 cm to facilitate the passage of waste and ensure proper function, while also allowing for effective management and care of the stoma.
\item \textbf{Answer:} False
\\ \textit{Options:} A) True; B) False
\\ \textit{Note:} The statement is false because while heterosexual males can be at risk, data show that gay and bisexual men, regardless of ethnicity, constitute the primary risk group for HIV infection, including among Latinos.
\item \textbf{Answer:} False
\\ \textit{Options:} A) True; B) False
\\ \textit{Note:} The frequency of AABBCC individuals from a mating of two AaBbCc individuals is calculated using a Punnett square and is actually 1/64, not 1/128, as there are 64 possible combinations of alleles in the offspring.
\end{enumerate}

\section*{Long Form Response}
\begin{enumerate}[label=\arabic*.]
\setcounter{enumi}{10}
\item \textbf{Answer:} Pharmacological agents can play a significant role in the etiology of gout by influencing uric acid metabolism and excretion. Certain diuretics, such as thiazides and loop diuretics, can lead to increased serum uric acid levels by promoting renal retention of uric acid, thereby increasing the likelihood of hyperuricemia and subsequent gout attacks. Additionally, medications like aspirin, particularly at low doses, may inhibit uric acid excretion, further contributing to the risk of gout. Other drug classes, including some immunosuppressants and anti-tuberculosis medications, have also been implicated in the induction of hyperuricemia. Understanding these mechanisms is crucial for healthcare professionals in both preventing and managing gout in patients who are on these pharmacological agents.
\\ \textit{Note:} correct because it identifies specific pharmacological agents, such as diuretics and low-dose aspirin, that interfere with uric acid metabolism and excretion, thereby increasing the risk of hyperuricemia and gout.
\item \textbf{Answer:} A congenital disorder that impairs leptin secretion leads to significant phenotypic consequences, primarily characterized by obesity and abnormal growth patterns. Leptin, a crucial hormone for regulating energy balance and appetite, when deficient, results in hyperphagia and decreased energy expenditure, contributing to obesity. Additionally, impaired leptin signaling can disrupt growth regulation, potentially leading to stunted growth or abnormal growth trajectories. This disorder is also associated with metabolic dysfunctions such as hypothyroidism and hyperinsulinaemia, where decreased leptin levels can impair thyroid hormone production and enhance insulin resistance, further exacerbating metabolic abnormalities. Collectively, these factors illustrate the critical role of leptin in maintaining metabolic homeostasis and normal growth.
\\ \textit{Note:} correct because it accurately describes how impaired leptin secretion disrupts energy balance and appetite regulation, leading to obesity and abnormal growth patterns, as leptin plays a critical role in both metabolic processes and growth regulation.
\end{enumerate}

\vfill
\noindent\textit{End of Answer Key}
\end{document}